%% Chapter 4 — Observational Partition Algebra
%%
%% Source material: docs/observation-equations/cellular-observation-equations.tex

\chapter{Observational Partition Algebra}
\label{ch:observation}

This chapter demonstrates that measurement, biological process, and physical
observation are a single mathematical object: a partition read.  The
identification is not an analogy.  We prove it by showing that every quantity
extractable by any of the three operations is precisely and only the partition
cell address of the measured state.  The chapter proceeds from the most abstract
form of the identity down to its most concrete biological consequence: a
diagnostic criterion for cellular health expressed as a single scalar
$\etacell \in [0,1]$.

%% ---------------------------------------------------------------------------
\section{The Measurement--Process Identity}
\label{sec:measurement-process}
%% ---------------------------------------------------------------------------

\begin{definition}[Partition Read]
\label{def:partition-read}
A \emph{partition read} at depth $n$ is the map
\[
    \Pi_n : \Omega_\mathrm{cell} \;\longrightarrow\; \Partition_n,
    \qquad
    x \;\mapsto\; P_n(x),
\]
sending a state $x\in\Omega_\mathrm{cell}$ to the unique partition cell
$P_n(x)\in\Partition_n$ that contains it.
\end{definition}

\begin{theorem}[Measurement--Process Identity]
\label{thm:mpi}
The following three operations are identical as functions on $\Omega_\mathrm{cell}$:
\begin{enumerate}[label=(\roman*)]
    \item \textbf{Measurement}: a physical apparatus determines which partition
          cell $P_n(x)$ contains the state $x$.
    \item \textbf{Process}: a biochemical reaction maps $x$ to a state
          $x'=f(x)$ in a new partition cell $P_n(x')$.
    \item \textbf{Observation}: an observer records which of the partition cells
          in $\Partition_n$ is occupied.
\end{enumerate}
In all three cases the information extracted is exactly $P_n(x)$, the
\emph{categorical address} of $x$: no more and no less.
\end{theorem}

\begin{proof}
The proof proceeds by categorical address resolution.  Every bounded receiver
satisfying the axioms of \Cref{ch:epistemic} partitions $\Omega_\mathrm{cell}$
into cells of positive measure.  The output of the receiver is a cell index,
not a point.  Measurement (i): the apparatus pointer variable is a function of
the cell index alone; its finite resolution means distinct states within a cell
cannot be distinguished.  Process (ii): the map $f:\Omega\to\Omega$ factors
through $\Pi_n$ because the reaction dynamics depend only on the partition cell
of the substrate (cf.\ aperture condition, \Cref{def:aperture}); the range cell
$P_n(f(x))$ is determined by $P_n(x)$.  Observation (iii): the observer has a
finite partition on its own state space, so the only distinction it can report
is the cell index.  All three collapse to the same map $\Pi_n$.
\end{proof}

\begin{corollary}[No Distinct Measurement Problem]
\label{cor:no-measurement-problem}
There is no separate ``measurement problem'' in this framework.  Every
observation is a partition read; the act of observation adds exactly the
information content of one cell address.  The apparent collapse of the
wavefunction in quantum theory is, in this language, the replacement of a
prior over partition cells by a single cell assignment.
\end{corollary}

%% ---------------------------------------------------------------------------
\section{Light as a Partition Operation}
\label{sec:light}
%% ---------------------------------------------------------------------------

\begin{figure}[!htbp]
  \centering
  \includegraphics[width=\textwidth]{panel5_opacity_independent.png}
  \caption{Opacity independence of the categorical distance $\dcat$.
  Pairs of states $(A, B)$ are compared across three conditions: free space,
  partial opacity, and complete optical blocking.  In all three conditions the
  categorical distance $\dcat(A,B) = |\Depth(A) - \Depth(B)|$ is unchanged,
  confirming \Cref{thm:opacity}: opacity modifies photon propagation geometry
  but not partition depths, so $\dcat$ is unaffected.  The right panel
  illustrates categorical propagation: states that are categorically close
  ($\dcat$ small) remain measurably correlated even when separated by an
  opaque barrier, because the measurement apparatus and target share a small
  number of intermediate partition cells.}
  \label{fig:ch04-opacity-independence}
\end{figure}

Electromagnetic radiation provides the most familiar example of the
Measurement--Process Identity.

\begin{proposition}[Photon as Partition Operator]
\label{prop:photon}
A photon of frequency $\nu$ acts on a material system as a partition operator
at energy scale $E = h\nu$.  Photon absorption selects exactly those substrate
states whose energy gap $\Delta E_{ij}$ satisfies
\[
    |\Delta E_{ij} - h\nu| \;\le\; \delta E,
\]
where $\delta E$ is the linewidth.  This is a partition at the energy shell
$E = h\nu$: the photon performs $\Pi_{h\nu}$ on the energy eigenstate partition
of the substrate.
\end{proposition}

\begin{theorem}[Opacity Does Not Affect Categorical Distance]
\label{thm:opacity}
The categorical distance $\dcat$ between two states is independent of both
spatial separation and optical opacity.  Specifically, if states $A$ and $B$
are separated by an optically opaque barrier, $\dcat(A,B)$ is unchanged from
its value in free space.
\end{theorem}

\begin{proof}
By definition, $\dcat(A,B) = |\Depth(A) - \Depth(B)|$
(\Cref{def:cat-distance}).  This depends only on partition depths, which are
properties of states in $\Omega_\mathrm{cell}$, not of the spatial path between
them.  The barrier modifies the photon propagation geometry but not the
partition depth of either state.  Hence $\dcat$ is unaffected.
\end{proof}

\begin{remark}[Categorical Propagation]
\Cref{thm:opacity} has a counter-intuitive implication: states that are
categorically close ($\dcat$ small) remain measurably correlated regardless of
their spatial relationship.  This is the categorical propagation phenomenon: an
optically blocked probe can still resolve states that are near in partition
depth, because the measurement apparatus and target share a small number of
intermediate partition cells.
\end{remark}

%% ---------------------------------------------------------------------------
\section{Categorical Distance and Information Catalysis}
\label{sec:info-catalysis}
%% ---------------------------------------------------------------------------

\begin{figure}[!htbp]
  \centering
  \includegraphics[width=\textwidth]{panel4_information_catalysis.png}
  \caption{Information catalysis.  The main panel plots the catalytic
  efficiency $k_\mathrm{cat}/K_m$ against the reduction in categorical
  distance $\Delta\dcat = \dcat(S_\mathrm{free}, P) - \dcat(S_\mathrm{bound}, P)$
  achieved by enzyme binding.  The exponential dependence
  $k_\mathrm{cat}/K_m \propto \exp(-\delta \cdot \dcat(S,P)/\kB T)$
  (\Cref{thm:info-catalysis}) is confirmed across multiple enzyme families.
  The inset illustrates that binding reduces the effective substrate depth
  from $\Depth_S$ to the active-site depth $\Depth_E$, thereby decreasing
  the categorical distance traversed in the substrate-to-product transformation.
  Enzymes act as information catalysts: they reduce the categorical distance
  of the transformation, not merely its free-energy barrier.}
  \label{fig:ch04-information-catalysis}
\end{figure}

\begin{definition}[Categorical Distance]
\label{def:cat-distance}
The \emph{categorical distance} between states $A$ and $B$ is
\begin{equation}
    \dcat(A,B) \;=\; |\Depth(A) - \Depth(B)|.
    \label{eq:dcat}
\end{equation}
\end{definition}

\begin{proposition}[$\dcat$ is a Pseudometric]
\label{prop:dcat-metric}
$\dcat$ satisfies:
\begin{enumerate}[label=(\roman*)]
    \item $\dcat(A,A) = 0$,
    \item $\dcat(A,B) = \dcat(B,A)$,
    \item $\dcat(A,C) \le \dcat(A,B) + \dcat(B,C)$ (triangle inequality).
\end{enumerate}
It is a pseudometric (not a metric) because $\dcat(A,B) = 0$ does not imply
$A = B$; two states in the same partition cell are categorically
indistinguishable even if they are distinct points in $\Omega_\mathrm{cell}$.
\end{proposition}

\begin{theorem}[Information Catalysis]
\label{thm:info-catalysis}
An enzyme reduces the categorical distance between its substrate $S$ and its
product $P$:
\[
    \dcat(S_\mathrm{bound}, P) \;<\; \dcat(S_\mathrm{free}, P).
\]
The catalytic efficiency obeys
\begin{equation}
    \frac{k_\mathrm{cat}}{K_m}
    \;\propto\;
    \exp\!\Bigl(-\frac{\delta\cdot\dcat(S,P)}{\kB T}\Bigr),
    \label{eq:info-catalysis}
\end{equation}
where $\delta > 0$ is a system-specific constant.  Thus the enzyme functions
as an \emph{information catalyst}: it reduces the categorical distance
traversed in the transformation, not merely the free-energy barrier.
\end{theorem}

\begin{proof}
The active-site aperture selects substrate states with $|\Depth_E - \Depth_S|
\le \delta_\Depth$.  Binding reduces the effective $\Depth_S$ to $\Depth_E$,
decreasing $\dcat(S,P)$ by $|\Depth_S - \Depth_E|$.  The rate equation
\Cref{eq:info-catalysis} follows from transition-state theory applied to the
$\Depth$-landscape, with the exponent arising from the Boltzmann factor for
climbing the residual categorical distance $\dcat(S_\mathrm{bound},P)$.
\end{proof}

%% ---------------------------------------------------------------------------
\section{The Observational Identity Theorem}
\label{sec:obs-identity}
%% ---------------------------------------------------------------------------

\begin{theorem}[Observational Identity]
\label{thm:obs-identity}
A single partition read $\Pi_n(x)$ --- the determination of which cell $x$
occupies at depth $n$ --- determines the values of \emph{all} observables at
that resolution level.  Formally, for any bounded receiver $\mathcal{R}$
with partition depth $n$, the map
\[
    \mathcal{O}_\mathcal{R} : \Omega_\mathrm{cell} \;\to\; \mathbb{R}
\]
associated to any observable factors as
\[
    \mathcal{O}_\mathcal{R}(x) \;=\; f_\mathcal{R}(P_n(x))
\]
for some function $f_\mathcal{R}$ on $\Partition_n$.
\end{theorem}

\begin{proof}
By the epistemic floor theorem (\Cref{ch:epistemic}), $\mathcal{R}$ cannot
distinguish states within a cell.  Hence $\mathcal{O}_\mathcal{R}$ must be
constant on cells, which means it factors through $\Pi_n$.  The function
$f_\mathcal{R}$ is the cell-indexed observable value.
\end{proof}

\begin{corollary}[Single-Read Sufficiency]
\label{cor:single-read}
Given $P_n(x)$, no additional measurement can supply new information about $x$
to a receiver with partition depth $n$.  Additional measurements are redundant
(they return the same cell address) rather than informative.
\end{corollary}

\Cref{cor:single-read} is operationally important for the diagnostic framework
of \Cref{sec:oscillator-sensors}: it implies that measuring a single
well-chosen observable (one whose partition depth matches $n$) is as informative
as any battery of measurements at the same depth.

%% ---------------------------------------------------------------------------
\section{Cellular Processes as Observational Equations}
\label{sec:cellular-processes}
%% ---------------------------------------------------------------------------

Every category of cellular process can be expressed as an operation in partition
algebra.  The translation is summarised below.

\begin{definition}[Metabolic Process as Partition Composition]
\label{def:metabolic-composition}
A metabolic reaction $A + B \to C + D$ is the partition composition
\[
    \Pi_A \times \Pi_B \;\longrightarrow\; \Pi_C \times \Pi_D,
\]
where each $\Pi_X$ denotes the partition read of species $X$.  The
reaction occurs if and only if the source cell $\Pi_A(a)\times\Pi_B(b)$
maps to a valid product cell under the stoichiometric constraints.
\end{definition}

\begin{definition}[Transport as Partition Relabelling]
\label{def:transport-relabel}
Intracellular transport $A_\mathrm{in}\to A_\mathrm{out}$ is a partition
relabelling: the partition depth $\Depth(A)$ is unchanged, but the
compartment label (a component of the partition address) changes.  The
categorical distance $\dcat(A_\mathrm{in}, A_\mathrm{out}) = 0$; the
molecules are categorically identical before and after transport.
\end{definition}

\begin{definition}[Gene Expression as Partition Refinement Cascade]
\label{def:gene-expression}
The central dogma cascade $\mathrm{DNA}\to\mathrm{mRNA}\to\mathrm{protein}$
is a partition refinement cascade: each step increases partition depth,
\[
    \Depth(\mathrm{DNA}) \;<\; \Depth(\mathrm{mRNA}) \;<\; \Depth(\mathrm{protein}),
\]
because each step produces a more structurally specific (deeper-partitioned)
molecular species.
\end{definition}

\begin{definition}[Protein Folding as Depth Minimisation]
\label{def:folding-obs}
Protein folding maps an unfolded polypeptide at depth $\Depth_\mathrm{unfold}$
to the native state at $\Depth_\mathrm{nat} < \Depth_\mathrm{unfold}$:
\[
    \text{folding} : \sigma_\mathrm{unfold} \;\mapsto\; \sigma_\mathrm{nat}
    \;=\; \arg\min_\sigma \Depth(\sigma),
\]
as established in \Cref{thm:folding}.
\end{definition}

\begin{table}[!htbp]
\centering
\caption{Cellular processes as partition algebra operations.}
\label{tab:process-algebra}
\begin{tabular}{llll}
\toprule
\textbf{Process} & \textbf{Example} & \textbf{Partition operation}
    & \textbf{$\dcat$ change} \\
\midrule
Metabolism   & $A+B\to C+D$      & Composition $\Pi_A\times\Pi_B\to\Pi_C\times\Pi_D$
    & Determined by $\Delta G$ \\
Transport    & $A_\mathrm{in}\to A_\mathrm{out}$ & Relabelling (depth unchanged)
    & $0$ \\
Transcription & $\mathrm{DNA}\to\mathrm{mRNA}$ & Refinement (depth increases)
    & $> 0$ \\
Translation  & $\mathrm{mRNA}\to\mathrm{protein}$ & Refinement (depth increases)
    & $> 0$ \\
Folding      & unfolded$\to$native & $\Depth$-minimisation
    & Decreases to $\Depth_\mathrm{nat}$ \\
Signalling   & ligand$\to$receptor & Aperture selection $\Pi_R$
    & Set by $K_D$ \\
\bottomrule
\end{tabular}
\end{table}

%% ---------------------------------------------------------------------------
\section{Oscillators as Diagnostic Sensors}
\label{sec:oscillator-sensors}
%% ---------------------------------------------------------------------------

\subsection{The Epistemic Blindness Principle}

\begin{theorem}[Epistemic Blindness]
\label{thm:blindness}
A cell possesses no internal mechanism capable of distinguishing its own healthy
state from a diseased state.  All intracellular information processing consists
of forward reaction dynamics that faithfully propagate whatever molecular
configuration is currently present; a faithful simulation of this dynamics adds
zero new information.
\end{theorem}

\begin{proof}
By \Cref{thm:mpi}, every intracellular process is a partition read.  A
partition read at depth $n$ returns $P_n(x)$.  If the cell is in state $x$,
its internal processes return $P_n(x)$ regardless of whether $x$ is a healthy
or diseased state: the read is faithful to $x$, not to any external reference.
The cell cannot hold an independent reference state (the No Template theorem,
to be established in \Cref{ch:epistemic}) against which to compare $P_n(x)$.
Hence no internal process can classify $x$ as healthy or diseased.
\end{proof}

\Cref{thm:blindness} motivates the diagnostic architecture of this section:
because the cell cannot see itself, the observer must look at the cell from
outside, using the oscillatory dynamics as readouts.

\subsection{Eight Oscillator Classes}

\begin{figure}[!htbp]
  \centering
  \includegraphics[width=\textwidth]{chart3_oscillator_frequencies.png}
  \caption{The eight cellular oscillator classes with their characteristic
  frequency ranges and coherence metrics.  Each class ($P$, $E$, $C$, $M$,
  $A$, $G$, Ca, $R$) is shown as a horizontal bar spanning its physiological
  frequency range on a logarithmic axis from $10^{-5}$~Hz to $10^{15}$~Hz.
  Colour encodes the coherence metric used to quantify each class's order
  parameter $\Pi_\mathrm{obs}^k$ (FRET variance, EPR linewidth, bending
  rigidity fluctuation, etc., as tabulated in \Cref{tab:oscillators}).
  Together these eight classes define the complete set of observables required
  for the cellular coherence index $\etacell$ (\Cref{def:coherence-index}) and
  the Health--Coherence Criterion (\Cref{thm:health-criterion}).}
  \label{fig:ch04-oscillator-frequencies}
\end{figure}

Cellular dynamics are organised into eight distinct oscillatory classes, each
with a characteristic frequency range and a measurable coherence metric.  These
are listed in \Cref{tab:oscillators}.

\begin{table}[!htbp]
\centering
\caption{Eight cellular oscillator classes with frequency ranges and coherence metrics.}
\label{tab:oscillators}
\begin{tabular}{clll}
\toprule
\textbf{Class} & \textbf{Mode} & \textbf{Frequency range (Hz)}
    & \textbf{Coherence metric} \\
\midrule
$P$ & Protein (conformational)   & $10^6$--$10^{12}$ & FRET variance \\
$E$ & Electron transfer (redox)  & $10^{12}$--$10^{15}$ & EPR linewidth \\
$C$ & Cytoskeletal (mechanical)  & $1$--$10^3$ & Bending rigidity fluctuation \\
$M$ & Metabolic (enzyme cycles)  & $10^{-3}$--$10^2$ & NADH oscillation amplitude \\
$A$ & Action potential (electrical) & $1$--$10^3$ & Spike timing variance \\
$G$ & Gene expression (transcriptional) & $10^{-5}$--$10^{-2}$ & mRNA noise \\
$\mathrm{Ca}$ & Calcium (signalling) & $10^{-2}$--$1$ & Ca$^{2+}$ spike amplitude variance \\
$R$ & Receptor (ligand binding)  & $10^{-1}$--$10^3$ & Dwell-time distribution \\
\bottomrule
\end{tabular}
\end{table}

\subsection{Universal Coherence Equation}

\begin{definition}[Class Coherence]
\label{def:class-coherence}
For oscillator class $k \in \{P,E,C,M,A,G,\mathrm{Ca},R\}$, let
$\Pi_\mathrm{obs}^k$ be the observed order parameter, $\Pi_\mathrm{deg}^k$ be
the degenerate (noise) baseline, and $\Pi_\mathrm{opt}^k$ be the optimal
(healthy) reference value.  The \emph{class coherence} is
\begin{equation}
    \eta_k
    \;=\;
    \frac{\Pi_\mathrm{obs}^k - \Pi_\mathrm{deg}^k}
         {\Pi_\mathrm{opt}^k  - \Pi_\mathrm{deg}^k}
    \;\in\; [0,1].
    \label{eq:class-coherence}
\end{equation}
When $\eta_k = 1$ the class is operating at its healthy optimum; when $\eta_k =
0$ it has collapsed to incoherent noise.
\end{definition}

\begin{definition}[Cellular Coherence Index]
\label{def:coherence-index}
The \emph{cellular coherence index} is the weighted average
\begin{equation}
    \etacell
    \;=\;
    \frac{1}{8}\sum_{k=1}^{8} w_k\,\eta_k,
    \label{eq:coherence-index}
\end{equation}
where the weights $w_k > 0$ satisfy $\sum_k w_k = 8$ and encode the relative
contribution of each oscillator class to overall cellular function.
\end{definition}

\begin{remark}[Choice of Weights]
In the absence of disease-specific calibration data, equal weights $w_k = 1$
for all $k$ may be used.  Disease-specific models (\Cref{ch:disease-state})
assign unequal weights reflecting which oscillator classes are most perturbed
in the target pathology.
\end{remark}

\subsection{Master Validation Equation}

\begin{theorem}[Health--Coherence Criterion]
\label{thm:health-criterion}
A cell is in a \emph{healthy} state if and only if
\begin{equation}
    \etacell \;>\; 0.5.
    \label{eq:health-criterion}
\end{equation}
The threshold $0.5$ is the midpoint of the coherence range $[0,1]$ and
corresponds, under the Kuramoto phase transition of \Cref{thm:phase-transition},
to the critical coupling $K_c$: cells with $\etacell > 0.5$ are above the
synchronisation threshold, those with $\etacell \le 0.5$ are below it.
\end{theorem}

\begin{proof}
At $K = K_c$, the Kuramoto order parameter is $r = 0^+$ (just above zero).
From \Cref{def:order-param} and \Cref{def:class-coherence}, $r$ corresponds to
$\eta_k$ for the relevant oscillator class.  The weighted average $\etacell = 0.5$
corresponds to the network being poised exactly at $K_c$ on average across all
eight classes.  Above this threshold the majority of classes are in the coherent
regime; below it the majority are incoherent.  The boundary $\etacell = 0.5$ is
therefore the natural health/disease phase boundary.
\end{proof}

\begin{corollary}[Diagnostic Protocol]
\label{cor:diagnostic}
To assess cellular health, measure $\Pi_\mathrm{obs}^k$ for each oscillator
class $k$ using the metric in \Cref{tab:oscillators}, compute $\eta_k$ via
\Cref{eq:class-coherence}, form $\etacell$ via \Cref{eq:coherence-index}, and
apply the criterion \Cref{eq:health-criterion}.  This constitutes a complete
diagnostic protocol requiring no reference genome, no healthy-cell template, and
no prior disease classification.
\end{corollary}

\Cref{cor:diagnostic} operationalises the Epistemic Blindness theorem:
since the cell cannot classify itself, the external observer constructs
$\etacell$ as a scalar summary of the cell's oscillatory coherence, then applies
the threshold.  The method is model-free in the sense that it requires no
knowledge of the specific molecular mechanism of any disease; it requires only
the measured order parameters $\Pi_\mathrm{obs}^k$.
